% =======================================================================
\chapter*{Resumen}
\addcontentsline{toc}{chapter}{Resumen}
% =======================================================================

Los Problemas de Optimización Combinatoria (POC) son mayoritariamente
\NP-duros; en la práctica se abordan con metaheurísticas. La
Búsqueda Tabú es una de las más establecidas, y admite enriquecerse con
mecanismos de memoria de largo plazo —lista tabú dinámica, memoria de
frecuencia, penalización adaptativa, diversificación informada— cuya
contribución individual rara vez se cuantifica: la literatura compara familias completas de
metaheurísticas o propone variantes ajustadas a un único problema, pero no
aísla qué componente interno hace el trabajo.

Este trabajo cierra esa brecha con un diseño factorial completo. Se
implementan tres configuraciones acumulativas de Búsqueda Tabú —Básica,
Intermedia y Avanzada— y se aplican, sin re-ajustes por instancia, a cinco
problemas de estructura heterogénea: Mochila 0/1, TSP, CVRP, Flowshop y Job
Shop. Sobre 24 instancias de benchmarks públicos con mejor valor conocido
documentado se ejecutan \textbf{360 experimentos} (24 instancias $\times$
3 niveles $\times$ 5 semillas independientes), y en cada uno se mide el GAP
porcentual respecto al BKS.

La transición de la configuración Básica a la Avanzada reduce el GAP global
promedio de \SI{6.50}{\percent} a \SI{3.40}{\percent}, una mejora del
\SI{48}{\percent} relativo, a un costo de aproximadamente $3.9\times$ el tiempo
de cómputo. La ganancia, sin embargo, no es uniforme: se concentra casi por
completo en CVRP (\SI{19.22}{\percent} $\rightarrow$ \SI{6.38}{\percent}) y es
marginal en Flowshop, donde la vecindad de intercambio adyacente apenas mejora
la solución de la heurística constructiva NEH y los tres niveles quedan en
torno al \SI{1.7}{\percent}. El hallazgo de mayor valor práctico apunta al primer salto de complejidad, donde se obtiene la mayor parte de la mejora: la configuración Intermedia captura 2.50 de los 3.10 puntos
porcentuales totales.

El producto principal es un repositorio público, reproducible y extensible que
contiene las quince implementaciones, el orquestador de experimentos, modelos
exactos de verificación y los resultados de las 360 corridas.

\vspace{1.5em}
\noindent\textbf{Palabras clave:} búsqueda tabú, CVRP, flowshop, job shop,
metaheurísticas, mochila 0/1, optimización combinatoria, repositorio público,
reproducibilidad, TSP.
